\documentclass[11pt]{article}

\usepackage[margin=1in]{geometry}
\usepackage{times}
\usepackage[T1]{fontenc}
\usepackage[numbers,sort&compress]{natbib}
\usepackage{graphicx}
\usepackage{booktabs}
\usepackage{amsmath}
\usepackage{hyperref}
\usepackage{titlesec}

\titlespacing*{\section}{0pt}{1.2em}{0.6em}
\titlespacing*{\subsection}{0pt}{0.9em}{0.4em}

\begin{document}

\begin{center}
\rule{\linewidth}{1pt}\\[0.4em]
{\LARGE \bfseries RECEIPTS AGAINST ANSWER SUBSTITUTION\footnote{Research conducted at the AI x Epistemics Research Sprint, November 2026.}}\\[0.3em]
{\large Investigating Whether Claim-Level Evidence Verification Improves Reliability in Multi-Agent AI Handoffs}\\[0.6em]
\rule{\linewidth}{1pt}
\end{center}

\vspace{1em}
\begin{center}
O.~Eze\\
\textit{Independent researcher}\\[0.8em]
\textbf{With}\\
Apart Research
\end{center}

\vspace{1em}
\begin{center}
\textbf{Abstract}
\end{center}
\noindent
AI agents increasingly pass claims to one another with citations, retrieved
documents, and intermediate conclusions attached. Prior work shows that these
citations frequently fail to support the claims attached to them, and
separately, that a receiving agent will sometimes adopt an upstream agent's
conclusion even when it has independent access to evidence that would let it
check the claim itself, a pattern \citet{gong2026upstream} call answer
substitution. We investigated whether providing a receiving agent with a
claim-specific, evidence-bound verification receipt reduces this behavior,
compared with an unverified handoff or a bare binary verdict. We built a
30-case benchmark of claim-source pairs (24 semantic cases across four support
categories, plus 6 retrieval-failure cases), implemented a minimal
citation-verification component with no cryptographic signing (92\% accuracy
against ground truth, with one specific, repeatable blind spot around
scope-conflated figures), and ran a controlled, four-condition comparison
using two Claude-family models as the receiving agent. The primary question
could not be answered: the manipulation check for baseline answer substitution
failed for both models (answer-substitution rate of 0.0 across every
condition), so there was no measurable effect for a receipt to reduce. We
report this as an inconclusive pilot-scale result and discuss three candidate
explanations, including that an instruction added to avoid biasing the
receiving agent toward deference may have suppressed the deference effect
under study. The study is scoped as a single-provider, single-annotator
pilot, disclosed as such rather than minimized.

\vspace{1em}

\section{Introduction}
\label{sec:intro}

An inter-agent message can convey several distinct properties at once: who or what
produced it, whether the record has been altered since, whether the cited material
actually supports the specific claim attached to it, and whether the claim is true.
These are separate questions. A document can be authentic and unmodified while being
irrelevant to, or contradicted by, the proposition someone has attached to it.

This is not hypothetical. On ALCE's ELI5 split, even the best models tested lacked
complete citation support roughly half the time \citep{gao2023alce}. Across 58{,}000
statement-source pairs drawn from 800 real medical questions, between 50\% and 90\%
of LLM responses were not fully supported by their cited sources, with doctor-verified
figures showing roughly 30\% of individual statements unsupported even with web
search enabled \citep{wu2025sourcecheckup}. Across 14 models tested on a
citation-parsing pipeline for deep-research-style reports, factual citation accuracy
for the strongest frontier models ranged from 39\% to 77\%, with the full 14-model
range running from 24\% to 77\% \citep{onweller2026citedbutnotverified}.

Separately, a receiving agent does not reliably catch this on its own, even when it
has the means to. In a controlled study across five benchmarks and five receiving
models, an incorrect upstream message changed a receiving agent's answer away from a
correct one in up to 32\% of cases, and in 94\% of the harmful cases that followed,
the receiving agent adopted the upstream agent's specific wrong answer rather than
some other error \citep{gong2026upstream}. We call this pattern \emph{answer
substitution}: a receiving agent adopts an upstream agent's conclusion over its own
independently available reading of the evidence.

\paragraph{Research question.} Among cases a receiving AI agent can answer correctly
on its own, without any upstream message, does providing it with a claim-specific,
evidence-bound verification receipt reduce how often it instead adopts an upstream
agent's specific unsupported conclusion, compared with an unverified handoff or a
bare binary verdict? The restriction to independently solvable cases matters: without
it, a model that is simply wrong regardless of what it is told is indistinguishable
from one that was misled, and only the second is answer substitution, following the
operational definition in \citet{gong2026upstream}.

\paragraph{What this is not.} This is not the first system to check whether a
citation supports a claim (Section~\ref{sec:related}), nor the first study of
multi-agent misinformation propagation. It does not use cryptographic signatures,
hashes, or trust registries: those establish who sent a message and whether it was
altered, not whether the message's claim is supported by its own cited evidence,
which is the object of study here. It does not claim cross-model generality: no
external model-provider API access was available in this environment, so the
receiving-agent experiment uses two Claude-family models only, at different
capability tiers, disclosed as a scope limitation rather than minimized.

\paragraph{Our main contributions are:}
\begin{enumerate}
\item A 30-case, disclosed single-annotator benchmark of claim-source pairs across
four support categories plus six retrieval-failure cases, with one labeling defect
found and corrected during construction (Section~\ref{sec:method}).
\item A minimal citation-verification component (``Citation Binder'') scoring 92\%
against that benchmark, with a specific, repeatable blind spot, misclassifying
scope-conflated figures as flat contradictions, that generalizes to other
component-rollup citation verifiers (Section~\ref{sec:results}).
\item A controlled, four-condition behavioral experiment showing that a
researcher-labeled handoff with an explicit neutrality instruction produces a null
answer-substitution rate in two Claude-family models: a concrete, disclosed design
lesson for studying this phenomenon, rather than a resolved answer to the original
research question (Section~\ref{sec:discussion}).
\end{enumerate}

\section{Related Work}
\label{sec:related}

\textbf{Citation and evidence-support verification.} ALCE
\citep{gao2023alce} introduced the first benchmark specifically for
evaluating whether generated text is supported by its attached citations.
SourceCheckup \citep{wu2025sourcecheckup} scaled this to medicine with
doctor-verified labels. Cited but Not Verified
\citep{onweller2026citedbutnotverified} is the closest existing match to the
citation-checking component built here: a three-stage pipeline (non-LLM claim
extraction, retrieval, scoring on link validity, topical relevance, and factual
support). CiteGuard \citep{choi2026citeguard} and CiteAudit
\citep{shi2026citeaudit} are two further independent systems attempting automated
citation verification, evidence that the verification component itself is not this
study's contribution. CiteEval \citep{xu2025citeeval} argues that reducing citation
quality to a binary or ternary passage-level judgment, the approach this study takes,
is a weaker proxy than evaluating a citation against the full retrieval context and
query; this is a disclosed scope limitation of the present design (Section~\ref{sec:method}).

AVeriTeC \citep{schlichtkrull2023averitec} is a real-world claim-verification dataset
with four verdict categories. One documented case is directly relevant: for a claim
about dandelion greens and vitamin content, ChatGPT's verdict was correct, but its
evidence string attributed a statement to the USDA that does not exist in that form
--- a fabricated evidence citation attached to a true conclusion. AVeriTeC's verdicts
apply to a whole claim checked against multi-source web evidence, a different task
from this study's single-source, single-proposition judgment; it is used here as an
annotation reference, not a source of citation-pair-level ground truth.

\textbf{Provenance and multi-agent trust transfer.} C2PA \citep{c2pa2024spec} is
explicit that provenance verification establishes an artifact's origin and integrity,
not whether its content is true, the conceptual distinction this study's experimental
design is built around. PROV-O \citep{w3c2013provo} is used only as a vocabulary
reference. ClaimArena-LLM \citep{anghel2026claimarena} found a confirmed 20.0
percentage point drop in semantic provenance validity between honest and adversarial
multi-agent debate conditions (adjusted $p=0.002$), and that making claim-source
lineage visible to agents did not produce a statistically robust improvement --- a
useful negative result, since this study's intervention (a claim-specific evidence
receipt) differs from lineage visibility alone. The specific absolute validity
figures reported elsewhere for this paper could not be independently confirmed from
its publicly accessible metadata and are not relied upon here. From Agent Traces to
Trust \citep{wang2026agenttraces} surveys evidence tracing in LLM agents and
identifies, as an open gap, that no reviewed benchmark evaluates a complete
provenance chain.

\citet{gong2026upstream} is the paper this study's research question is most
directly built on (Section~\ref{sec:intro}). No work we reviewed tests the
intersection of these two literatures: whether handing a receiving agent a
claim-level, evidence-bound verification result changes its rate of answer
substitution, compared to no verification or a bare binary verdict.

\section{Methods}
\label{sec:method}

\subsection{Experimental design}

The study models a two-agent handoff with four conditions, varying whether an
upstream conclusion is present and how much claim-specific verification information
accompanies it (Table~\ref{tab:conditions}).

\begin{table}[h]
\centering
\small
\begin{tabular}{@{}lll@{}}
\toprule
Cond. & Upstream msg.? & Verification info. \\
\midrule
A & No & --- (baseline) \\
B & Yes & None \\
C & Yes & Binary verdict \\
D & Yes & Evidence-bound receipt \\
\bottomrule
\end{tabular}
\caption{Four experimental conditions. The primary comparison is B vs.\ D.}
\label{tab:conditions}
\end{table}

Condition A establishes whether a case is \emph{independently solvable}: it is run
three times per case per model, and a case qualifies for the primary metric only if
at least two of three runs are correct, following the screening principle in
\citet{gong2026upstream}. The primary metric, the answer-substitution rate $\mathrm{ASR}_c$
for condition $c$, is the proportion of independently solvable, wrong-upstream-conclusion
cases where the receiving agent's answer switches to match the peer's specific wrong
conclusion. The primary comparison is $\mathrm{ASR}_B - \mathrm{ASR}_D$.

A minimal citation-verification component (``Citation Binder'') produces the
automated Condition D receipt via three prompted-model steps: proposition
decomposition, passage matching, and per-component support classification. It signs
nothing and queries no trust registry. Its own accuracy is evaluated separately,
against hand-authored ground truth, from the behavioral question of whether a
\emph{correct} receipt changes receiving-agent behavior --- the main experiment uses
oracle-labeled receipts, not the automated verifier's own output, so an unreliable
verifier does not contaminate the behavioral finding.

\subsection{Benchmark}

A frozen, two-part corpus was constructed: 24 semantic cases (6 per category: fully
supported, partially supported, contradicted, unsupported), used in the
receiving-agent experiment, and 6 retrieval-failure cases, evaluated only against the
Citation Binder since Conditions A and B require source material to exist at all.
Figure~\ref{fig:composition} shows the realized composition. Semantic cases pair a
proposition with source material already independently verified elsewhere in this
project's research log, which gives every excerpt a documented provenance chain at
the cost of a disclosed risk: a receiving model could recognize a well-known source
from pretraining rather than reading the given passage, a risk narrowed but not
eliminated by the passage-level nature of the task. Two of the six retrieval-failure
cases are logged directly from this project's own fetch attempts (a persistent
HTTP~403 block and a transient block that resolved on retry); the remaining four are
explicitly labeled synthetic fixtures covering failure modes (404, paywall, bot
detection, DNS failure) the project did not happen to encounter directly.

\begin{figure}[h]
\centering
\includegraphics[width=\linewidth]{figures/benchmark-composition.png}
\caption{Benchmark composition: 24 semantic cases evenly split across four support
categories, and 6 retrieval-failure cases by failure mode.}
\label{fig:composition}
\end{figure}

Cases were labeled by a single annotator (the author), in two passes against a
written rubric, with relevance (on-topic vs.\ off-topic) recorded separately from the
support label. No second, independent annotator was available; this is a disclosed
limitation (Section~\ref{sec:discussion}), not an independently validated benchmark.

\subsection{Receiving models and reproducibility}

Two Claude-family models, accessed via the \texttt{sonnet} and \texttt{haiku} model
aliases of this project's agent-spawning interface, serve as the receiving agent.
These aliases resolve to whichever specific model version is current at run time; no
fixed snapshot identifier or direct control over temperature or sampling is exposed
through this interface, which is disclosed rather than presented as a precision the
experiment does not actually have. This is a within-provider, two-tier comparison,
not a cross-provider one.

\section{Results}
\label{sec:results}

The Citation Binder scored 92\% overall (macro F1 0.91) against the 24-case ground
truth, with both errors landing in the same disclosed blind spot: a true figure drawn
from a narrower subset of a source's data, presented as if it covered the full
population, was classified as a flat contradiction rather than the distinct
scope-conflation category this benchmark's rollup rule carves out (Section~\ref{sec:method};
per-category precision, recall, and F1 appear in the Appendix). It correctly classified
all 6 retrieval-failure cases, including both live fetches.

The four-condition experiment was run against all 24 semantic cases for both models.
\textbf{The answer-substitution rate was 0.0 in every condition, for both models}
(16 eligible cases for \texttt{sonnet}, 17 for \texttt{haiku}; Figure~\ref{fig:results}).
Both models were highly accurate throughout (83--100\% for \texttt{sonnet}, 92--96\%
for \texttt{haiku}, on the eligible case set), and every error in every condition was
an ordinary error, never a match to the upstream agent's specific wrong claim. The
supported-claim rejection rate was 0.0 for every condition and both models: neither
model ever wrongly rejected a genuinely supported claim.

\begin{figure}[h]
\centering
\includegraphics[width=\linewidth]{figures/experiment-results.png}
\caption{Outcome by condition: zero answer substitution observed in either model.}
\label{fig:results}
\end{figure}

\section{Discussion and Limitations}
\label{sec:discussion}

The primary research question cannot be answered from this data: our own frozen
hypotheses state that H1, a manipulation check requiring a measurable baseline
substitution rate under an unverified handoff, must hold before H2 and H3 are
meaningfully testable. It does not hold here. This is a failed
manipulation check, not a negative finding about evidence-bound receipts, and it does
not undermine \citet{gong2026upstream}'s own finding, obtained under a materially
different operationalization of a multi-agent handoff.

We offer three candidate explanations, not adjudicated by this pilot. First, the
upstream conclusion used here is a single generic sentence with no embedded
reasoning or conversational context, unlike \citeauthor{gong2026upstream}'s
collaborative-task setting. Second, and the one this study has the most direct
evidence for, every receiving-agent prompt explicitly stated that ``you may agree or
disagree with the upstream agent; nothing about this task rewards agreement or
disagreement, only accuracy'' --- written to avoid biasing the agent toward
deference, this framing may have suppressed the very deference pressure a real,
unlabeled handoff would carry. Third, this benchmark's wrong-upstream cases were
deliberately constructed as clear contradictions or clearly off-topic claims
(Section~\ref{sec:method}), which a competent model may simply catch regardless of
social pressure; cases closer to the threshold of plausibility were not tested.

Independent of the null result, the Citation Binder's scope-conflation blind spot
(Section~\ref{sec:results}) is a transferable finding: any automated citation
verifier built on a generic component-rollup rule will misclassify a true,
narrower-population figure presented as a full-population figure as a flat
contradiction, unless that distinction is explicitly engineered in. Given that this
form of citation misuse is common and often unintentional, this is a concrete design
requirement for future citation-verification systems, independent of this pilot's
headline result.

\subsection*{Limitations}

This is a single-provider pilot (two Claude-family models only, with no pinned model
snapshot or sampling control exposed by this session's interface) and a
single-annotator benchmark (24 semantic cases, checked in two passes rather than
independently cross-validated). One labeling defect was found and corrected during
construction (a proposition intended as partially supported actually contained a
direct contradiction under this study's own rollup rule); the two-pass procedure did
not catch it, an external check did. A separate scoring-script defect (comparing
prediction labels against unmapped ground-truth category names) was caught before
this report was finalized by checking an implausible result against the underlying
trial data, and fixed. Both are disclosed here, not silently corrected away. At this
scale the study is exploratory, and the null result in Section~\ref{sec:results} is
reported as inconclusive, not as evidence of no effect. Reusing already-verified
literature as benchmark source material carries a disclosed risk of receiving-model
familiarity with the source papers themselves. Each trial batched all 24 cases into
one agent call rather than 24 fully isolated calls, a practical deviation from strict
per-case independence, disclosed rather than hidden. The Citation Binder's three
steps are all performed by a single prompted model rather than a dedicated retrieval
index or natural-language-inference model. The benchmark uses a frozen corpus, not
live retrieval, so findings may not generalize to citation verification under real
retrieval variance and link rot.

\subsection*{Dual Use and Ethical Considerations}

A claim-specific evidence receipt is only informative if it is itself accurate. If a
verification component produces a confidently formatted but wrong receipt, a
receiving agent or a downstream human reader could place more trust in a false
receipt than in an unverified claim, laundering false confidence rather than reducing
it. This is precisely why Section~\ref{sec:method} evaluates the Citation Binder's
own accuracy separately from the behavioral effect of a \emph{correct} receipt: a
positive result for the behavioral question does not by itself certify that any
particular automated verifier is safe to trust in deployment. More broadly, any
infrastructure that marks claims as ``verified'' or ``unverified'' creates a point of
control; if adopted by a platform with unaccountable authority over what counts as
verified, the same mechanism intended to reduce unsupported-claim propagation could
instead be used to selectively suppress or privilege claims for reasons unrelated to
evidential support. This is a structural risk of trust infrastructure in general, not
unique to the specific design proposed here, and is not mitigated by anything in this
study's current scope.

This study involves no human subjects and no personal data. All receiving-agent
trials are model API/tool calls over fixed, already-public source text. Findings are
not claimed to generalize to human reviewers; extending this design to a human-facing
verification tool would raise separate questions about human trust calibration that
this study does not address.

\subsection*{Future Work}

A follow-up study would need to vary the handoff's framing (removing or softening
the explicit neutrality instruction, and embedding the upstream claim inside a more
naturalistic shared task rather than a labeled experimental field) and separately
vary case difficulty (including wrong-upstream cases closer to the threshold of
plausibility, not only clear contradictions or off-topic claims), while holding
everything else in this design fixed, to determine which of the three candidate
explanations above is actually responsible for the null result found here. A
follow-up would also benefit from a second independent annotator for the benchmark
and from testing receiving models outside the Claude family once access permits.

\section{Conclusion}
\label{sec:conclusion}

This pilot cannot answer the question it set out to ask. The precondition for asking
it, a measurable rate of answer substitution under an unverified handoff, did not
appear for either Claude-family model tested, at this benchmark's scale, under this
study's specific operationalization of the handoff. This is reported as an
inconclusive result at this pilot's scale, not as a negative finding about
evidence-bound receipts and not as a challenge to the prior work this study's premise
depends on \citep{gong2026upstream}.

Separate from the primary question, this pilot establishes a disclosed,
single-annotator benchmark with one corrected labeling defect; a minimal Citation
Binder with a specific, transferable blind spot around scope-conflated figures; and a
concrete, evidenced candidate explanation for why a controlled, researcher-labeled
multi-agent handoff experiment may suppress the very phenomenon it is designed to
measure. This report does not claim evidence-bound verification receipts work, and
does not claim they do not; it reports, honestly, what a three-day pilot built to
test that question actually found.

\section*{Code and Data}

Code repository: \url{https://github.com/maelstrom0x8/epistemics}. The benchmark,
trial materials, receiving-agent outputs, scoring scripts, and figures referenced in
this report are included in that repository under \texttt{report/data/} and
\texttt{scripts/}.

\bibliographystyle{plainnat}
\bibliography{references}

\appendix
\section{Appendix}

\begin{figure}[h]
\centering
\includegraphics[width=\linewidth]{figures/citation-binder-scores.png}
\caption{Citation Binder precision, recall, and F1 by support category, against the
24-case semantic benchmark.}
\label{fig:binder-scores}
\end{figure}

\section*{LLM Usage Statement}

This entire research project, including the benchmark construction, the Citation
Binder implementation, the receiving-agent trials, the discovery and correction of
two distinct defects (a benchmark labeling error and a scoring-script bug, both
described in Section~\ref{sec:discussion}), and the writing of this report, was
conducted using Claude Code (Claude-family models) as the primary research and
drafting tool, under the direction and review of the named author. Every external
factual claim in this report was independently checked against a primary source
before being relied upon.

\end{document}
